\chapter*{Conclusions}
\addcontentsline{toc}{chapter}{Conclusions}
\markboth{Conclusions}{Conclusions}

This thesis presented the complete design, implementation, and experimental validation of \textit{AI Auto Doctor}---a cross-platform mobile diagnostic application implementing a two-stage analysis pipeline that combines deterministic Rule-Based sensor verification with generative AI explanation of vehicle fault codes.
The following conclusions directly and explicitly address each research task defined in the Introduction.

\section*{Hypothesis Verification}
\addcontentsline{toc}{section}{Hypothesis Verification}

The research hypothesis is confirmed: the two-stage Rule-Based~+~AI pipeline achieves statistically significant improvements over direct AI-only output across all measured dimensions.

\begin{itemize}
  \item \textbf{Cause identification accuracy:} increased from 61\,\% (AI-only) to 87\,\% (Rule+AI), a~$+$26 percentage-point improvement measured across three DTC scenarios (P0171, P0300, P0420), each evaluated in three repeated runs by five automotive domain experts.
  \item \textbf{Result consistency:} increased from 58\,\% to 94\,\% ($+$36\,pp), demonstrating that the deterministic Rule-Based layer eliminates the response variability inherent in stochastic generative AI outputs.
  \item \textbf{User explainability score:} improved from 3.2/5 to 4.5/5 ($+$1.3), confirming that sensor-grounded explanations are perceived as more credible and actionable than training-set-based estimates.
  \item \textbf{``Can Drive?'' safety verdict:} absent in all AI-only outputs, present and deterministic in all Rule+AI outputs, addressing a critical safety information gap that stochastic AI cannot reliably resolve.
\end{itemize}

\section*{Achievement of Research Tasks}
\addcontentsline{toc}{section}{Achievement of Research Tasks}

\begin{enumerate}

  \item \textbf{Task~1 --- Literature Review:}
  A systematic survey of twelve mobile OBD-II diagnostic applications (Chapter~\ref{ch:litreview}, Table~\ref{tab:sota}) identified the research gap: no existing solution cross-validates DTC codes with live sensor readings before generating an explanation.
  Three diagnostic AI approaches were compared (Table~\ref{tab:aicomparison}); the Rule-Based~+~LLM combination uniquely satisfies all six evaluation criteria simultaneously.
  Industry evidence from Cooke~et~al.~\cite{cooke2022} and Bouchard~et~al.~\cite{bouchard2021} quantified the real-world impact of the identified gap.

  \item \textbf{Task~2 --- Architecture and Design:}
  A three-layer architecture with 16 components (Table~\ref{tab:components}) was designed and specified.
  The \texttt{DiagnosticEngine} covers five DTC rule groups (Table~\ref{tab:rule_groups}) with a formal six-component confidence scoring algorithm (Equation~\ref{eq:confidence}, Table~\ref{tab:confidence}, maximum~95\,\%).
  The \texttt{HealthScoreCalculator} applies a 14-condition penalty schedule (Table~\ref{tab:healthscore}, Equation~\ref{eq:healthscore}).
  The BLE OBD-II communication stack supports three ELM327 UUID configurations.
  Gemini~1.5~Flash integration uses temperature~0.4 and structured JSON output with an offline fallback mechanism.

  \item \textbf{Task~3 --- Implementation and Testing:}
  Eight application screens were implemented, covering the complete user journey from BLE connection through real-time monitoring to two-stage diagnosis and session history.
  All ten functional requirements passed UAT on Android~14 and iOS~17 (Table~\ref{tab:uat}).
  Dashboard cold start measured at 2.7\,s; \texttt{DiagnosticEngine} execution under 2\,ms; all performance thresholds met (Table~\ref{tab:performance}).
  STRIDE security analysis identified and mitigated six threat categories (Table~\ref{tab:stride}).

  \item \textbf{Task~4 --- Comparative Experiment:}
  A controlled experiment with three DTC scenarios, three repeated runs each, and five expert evaluators confirmed the hypothesis across four quantitative metrics (Table~\ref{tab:experiment}).
  The Rule+AI approach outperforms AI-Only on every measured dimension while reducing average API response time by 0.3\,s through prompt optimisation.

  \item \textbf{Task~5 --- Conclusions:}
  The experimental results confirm that introducing a deterministic, sensor-validated rule layer before LLM invocation significantly reduces AI hallucinations and response variability in technical diagnostic systems.
  This principle is generalisable to any IoT or cyber-physical system domain where AI is applied to interpret multi-sensor data.

\end{enumerate}

\section*{Scientific and Practical Significance}
\addcontentsline{toc}{section}{Scientific and Practical Significance}

The primary scientific contribution of this thesis is the demonstration that routing a generative AI model through a deterministic, sensor-validated rule layer eliminates a substantial proportion of diagnostic hallucinations and reduces result variability.
The two-stage constraint mechanism proposed here is replicable and applicable to analogous domains---industrial equipment diagnostics, medical sensor monitoring, building management systems---where AI systems interpret complex multi-sensor data without guaranteed access to real-time measurements.

The practical significance is threefold.
For vehicle owners: any driver with a \$10--30 ELM327 adapter obtains a sensor-verified, explainable, actionable diagnosis without technical expertise.
For automotive safety: the deterministic ``Can Drive?'' verdict, grounded in measured overheating state and fault severity, addresses a safety gap that AI-only systems cannot reliably close.
For the automotive service industry: the quantified reduction in misdiagnosis from 39\,\% to 13\,\% (inverse of cause accuracy) has direct economic implications, reducing unnecessary repairs estimated at 200--2\,000~EUR per incident~\cite{cooke2022}.

\section*{Limitations and Future Work}
\addcontentsline{toc}{section}{Limitations and Future Work}

Three principal limitations affect the current implementation.
First, \texttt{RealObdSource} polls only three of the fourteen defined LiveData parameters (RPM, Speed, Coolant Temperature) in real-hardware mode; the remaining eleven parameters are stub values.
This reduces confidence scores $\conf$ below their theoretically achievable maximum and is the highest-priority next development task.
Second, the Mode~03 multi-packet DTC response parser is not implemented for real-hardware mode, limiting production DTC reading to the demo simulation.
Third, the application has not been validated on physical ELM327 adapters across a diverse set of real vehicles; all testing was conducted in demo mode.

Future work directions, in order of priority:
(1)~complete the 14-PID polling implementation in \texttt{RealObdSource} (STFT PID~06, LTFT PID~07, MAF PID~10, O2 PID~14, and others);
(2)~implement the Mode~03 multi-frame DTC response parser;
(3)~add VIN reading (Mode~09, PID~02) for automatic vehicle profile population and per-model rule threshold calibration;
(4)~expand \texttt{DiagnosticEngine} from five to fifteen DTC rule groups (transmission P07xx, EVAP P04xx, EGR P04xx, ABS/Traction C1xxx);
(5)~validate the system on a diverse fleet of real vehicles to quantify the effect of per-manufacturer sensor threshold variations on rule accuracy.

\clearpage
